\section{Getting started}

This section introduces the system requirements and installation steps for \pppx{}
on the supported platforms (Linux, macOS, and Windows).

%To download the latest version of \pppx{}, please have a look on \href{https://github.com/pppx-dev/pppx/releases}{this page}. Note that only the dynamically linked version of \pppx{} is provided for macOS and Windows, but a statically linked version called \texttt{pppx\_static} is additionally provided for Linux.



\subsection{System requirements}

Generally, there is no specific hardware requirement for \pppx{}.
To process a 24-hour GNSS observation file with a sampling rate of 30 s,
\pppx{} consumes around 150 MB of resident memory and one CPU core.
The resident memory consumed may, however, vary with the number of GNSS systems used
and the data sampling rate.

Note that \pppx{} only supports Macs with Apple silicon (i.e., ARM architecture)
and a macOS version newer than Monterey.



\subsection{Installation}

\subsubsection{Linux}
\label{sec:installation_linux}

This application is distributed as an \href{https://appimage.org/}{AppImage},
making it portable across most modern Linux distributions.
Specifically, \pppx{} is built on Ubuntu 20.04 and requires the host system
to have GNU C Library (\texttt{glibc}) version 2.31 or newer.

It is compatible with the following distributions (and any of their derivatives):

\begin{table}[h]
\centering
\renewcommand{\arraystretch}{1.25}
\begin{tabular}{p{0.35\textwidth} p{0.55\textwidth}}
    \hline
    Distribution Family & Supported Versions \\ 
    \hline
    Ubuntu & 20.04 LTS, 22.04 LTS, 24.04 LTS (and newer) \\ 
    Linux Mint & 20, 21, 22 (and newer) \\ 
    Debian & 11 ``Bullseye'', 12 ``Bookworm'' (and newer) \\ 
    Fedora & 32 (and newer) \\ 
    RHEL / Rocky / Alma & 9.x (and newer) \newline \textit{Note: RHEL 8 is not supported} \\ 
    openSUSE & Leap 15.3+, Tumbleweed \\ 
    Arch Linux / Manjaro & Rolling releases (Always updated) \\ 
    \hline
\end{tabular}
\end{table}

\vspace{1em}
\textit{\textbf{Not sure if your system is supported?}}

You can easily check your system's \texttt{glibc} version
by running the \texttt{ldd --version} command in your terminal.


Clone the repository and run the installer from its root folder:

\begin{minted}[frame=leftline]{bash}
git clone --depth 1 https://github.com/pppx-dev/pppx.git
cd pppx
./install.sh
\end{minted}

The installer downloads the latest \pppx{} executable, installs both \pppx{} and the
wrapper script \texttt{pppx.sh} to \verb|~/.local/bin|, and configures \texttt{pppx.sh}
to use the \texttt{table} files in the cloned repository. Keeping the repository is
recommended, as it provides the \texttt{table} files and examples used by \pppx{}.



\subsubsection{Windows}
\label{sec:installation_windows}

The easiest way to run \pppx{} on Windows is to use the Windows Subsystem for Linux (WSL);
simply follow the instructions in Section~\ref{sec:installation_linux}.
However, \pppx{} can also run natively on Windows
using PowerShell or Command Prompt (\texttt{cmd.exe}).


Clone the repository, then run the installer from its root folder in PowerShell.
It will download the executable, fetch the required DLLs, and add the software
(\path{C:\Users\<YOUR_USER_NAME>\AppData\Local\Programs\PPPx}) to your user \texttt{PATH}:

\begin{minted}[frame=leftline]{powershell}
git clone --depth 1 https://github.com/pppx-dev/pppx.git
cd pppx
.\install.ps1
\end{minted}

\textbf{Note}: You must restart PowerShell, cmd.exe, or Git Bash after the installation.

The installer also installs and configures the wrapper script \texttt{pppx.sh}
(pointing it at the \texttt{table} folder in the cloned repository). While \texttt{pppx.sh}
cannot run in PowerShell or \texttt{cmd.exe}, you can run it from \textbf{Git Bash}, which
is bundled with \href{https://git-scm.com/download/win}{Git for Windows} (the same tool
used for \texttt{git clone}). Product downloads rely on \texttt{curl}, which Git Bash
provides:

\begin{minted}[frame=leftline]{bash}
# In Git Bash
pppx.sh -c pppx.ini rinex/ZIM200CHE_R_20221000000_01D_30S_MO.rnx
\end{minted}


% \paragraph{With MSYS2}

% It is recommended to use \pppx{} on Windows 10 or higher.

% % For Windows 7 or 8 users, please refer to \href{https://www.msys2.org/docs/windows_support/}{this page} to install a old version of \texttt{MSYS2} for software installation.

% If you don't have \texttt{MSYS2} installed on your PC yet, please install it according to \href{https://www.msys2.org}{this page}.

% Then, open \texttt{MSYS2 MINGW64} and run the following comands:
% \begin{minted}[frame=leftline]{bash}
%   $ pacman -Syu   # update package library
%   $ pacman -S git mingw-w64-x86_64-ceres-solver

%   $ git clone git@github.com:pppx-dev/pppx.git
%   $ mkdir -p ${HOME}/.local/bin/
%   $ cp pppx.exe ${HOME}/.local/bin/
% \end{minted}



\subsubsection{macOS}
\label{sec:installation_macos}

\pppx{} only supports Macs with Apple silicon and a macOS version newer than Monterey.


\paragraph{Prerequisite}

If you do not have \href{https://brew.sh/}{Homebrew} on your Mac, please install it first:

\begin{minted}[breaklines, frame=leftline]{bash}
  $ /bin/bash -c "$(curl -fsSL https://raw.githubusercontent.com/Homebrew/install/HEAD/install.sh)"
\end{minted} 

\pppx{} is linked against the latest \texttt{ceres-solver} library, so make sure you have the
latest version installed as well in order to run \pppx{} properly on your Mac.
Run the following command in the Terminal application:

\begin{minted}[frame=leftline]{bash}
brew install ceres-solver
\end{minted}


Clone the repository and run the installer from its root folder in the Terminal:

\begin{minted}[frame=leftline]{bash}
git clone --depth 1 https://github.com/pppx-dev/pppx.git
cd pppx
./install.sh
\end{minted}

The installer downloads the latest \pppx{} executable, bypasses Apple's quarantine
security check, installs both \pppx{} and the wrapper script \texttt{pppx.sh} to
\verb|~/.local/bin|, and configures \texttt{pppx.sh} to use the \texttt{table} files in
the cloned repository. Keeping the repository is recommended, as it provides the
\texttt{table} files and examples used by \pppx{}.
